%% notes-tex/database/database-expansion-200/sections/state.tex
%% [[experiments.database.expansion-200]]
%% https://github.com/Mjvolk3/torchcell/tree/main/notes-tex/database/database-expansion-200/sections/state.tex
%% SOURCE: hand-authored prose; record counts from database/releases/2026.10.17-4295f7ca.json,
%% row figures from experiments/database/scripts/build_candidate_datasets_table.py, and the
%% per-isolate gene-content figures from
%% experiments/database/scripts/isolate_gene_content_vs_s288c.py

\section{State across genomes\secstatus{todo}}\label{sec:state}

The \term{state} of a cell, for the purposes of this list, is its transcriptome or its
proteome: the abundance of every transcript or every protein at once. A \term{genome} is a
strain's total genomic content, written down as sequence. The first band of this ranking
is the set of rows that measure a state on more than one genome, where the genomes differ
by sequence rather than by one designed edit: natural isolates, the segregants of a cross,
a recombinant panel, or an interspecies hybrid read allele by allele. The band comes first
because a model that predicts state from genome has to be fitted on pairs of the two, and
the served graph holds far more of each half than of the pair.

\subsection{What the served graph holds on that axis\secstatus{todo}}\label{sec:state-served}

Table~\ref{tab:state-served} is every served dataset whose genotype axis is a set of
sequenced genomes rather than the S288C reference with one cataloged edit. Three of the
ten measure a state. Caudal 2024 is 943 isolate transcriptomes on the 1,011-isolate panel;
Muenzner 2024 is 795 proteomes on the same panel and 55 turnover profiles; Boocock 2025
and N'Guessan 2025 are single-cell transcriptomes over 82,467 segregant genotypes of four
crosses. Everything else on a sequenced-genome axis is growth, a life-cycle trait, or a
guide-by-background fitness, which fixes the genotype side of the pair without supplying a
state.

\begin{table}[H]\centering
\small
\caption[]{Served datasets whose genotype axis is a set of sequenced genomes. Records are
the served counts of release 2026.10.17-4295f7ca. \emph{State} marks a transcriptome or
proteome readout.}
\label{tab:state-served}
\begin{tabular}{@{}l l l r c@{}}
\toprule
Dataset & Genomes & Readout & Records & State \\
\midrule
Caudal 2024 & 943 natural isolates & RNA-seq TPM per gene & 943 & yes \\
Muenzner 2024 & 795 natural isolates & protein abundance per protein & 795 & yes \\
Muenzner 2024 & 55 natural isolates & protein degradation rate & 55 & yes \\
Boocock 2025 & 79,516 segregants, 3 crosses & single-cell UMI counts & 79,516 & yes \\
N'Guessan 2025 & 2,951 BY x RM segregants & single-cell UMI counts, pseudobulk & 2,951 & yes \\
Hale 2024 & 169 segregants x 8,046 guides & CRISPRi fitness & 8,692,635 & \\
Bloom 2019 & 13,950 segregants, 16 crosses & growth over 38 traits & 530,100 & \\
Peter 2018 & 971 natural isolates & growth over 35 conditions & 33,985 & \\
De Chiara 2022 & 1,011-isolate panel & doubling time, cell yield, life cycle & 76,620 & \\
Puddu 2019 & 4,800 deletion strains, sequenced & genome architecture & 8,835 & \\
\bottomrule
\end{tabular}
\end{table}

How different those genomes are from the reference is a measured quantity, not a
figure of speech. Against S288C, the 943 isolates of the served Caudal 2024 store carry a
median of 78 accessory open reading frames the reference lacks and lack a median of 126
the reference has (ranges 34 to 352 and 43 to 453), so a typical isolate is about 50 ORFs
smaller than BY4741 and differs from it at about 200 ORFs in presence. The larger
difference is inside the shared genes: a median of about 5,000 sequence variants and
roughly 360 natural loss-of-function alleles per isolate. Across the panel, 1,614 distinct
non-reference ORFs appear in at least one isolate and 1,118 reference ORFs are absent in at
least one. The population catalog these sit in is Peter 2018's pangenome of 7,796 ORFs,
4,940 core and 2,856 variable.

\subsection{What the eight rows add\secstatus{todo}}\label{sec:state-rows}

The eight rows of the band, ranks 1 to 8 of Table~\ref{tab:final}, each put a state on a
genome axis the graph does not yet pair it with.

\notesec{Albert 2018 and Jakobson 2025, the segregant pair} Albert 2018 is bulk RNA-seq of
1,012 sequenced BY x RM segregants, the largest genotype-indexed transcriptome in yeast and
the bulk partner of the served Boocock 2025, whose pooled panel holds 393 of these same
segregants read per cell. Jakobson 2025 measures protein abundance on the same sequenced
segregant class, 800 F6 segregants with more than 6,400 protein quantitative trait loci.
Together they are transcript and protein on one recombinant panel, which is the
natural-variation counterpart of the served Kemmeren 2014 and Messner 2023 pair on
deletions.

\notesec{Teyssonniere 2024 and Grossbach 2022, within-strain pairing} Teyssonniere 2024 is
proteomes for 942 of the 1,011 isolates, so against the served Caudal 2024 the
transcript-to-protein pairing is one to one by strain at species scale. It shares authors
and a panel with the served Muenzner 2024, and whether the two proteomes are independent
acquisitions is the first thing to settle, since building both without settling it risks
serving one measurement twice. Grossbach 2022 is the one row where transcript, protein and
phosphosite are all observed within a genotype, on 112 sequenced BY x RM strains, which
turns the RNA-to-protein question from a population correlation into a within-strain
residual.

\notesec{Smith 2008, the condition axis} Smith 2008 profiles 109 BY x RM segregants in
glucose and in ethanol. It is the only segregant transcriptome with an environment axis,
so genotype by environment on expression is a contrast inside one panel rather than a join
across two.

\notesec{Skelly 2013, McManus 2014 and Foss 2007, the smaller rows} Skelly 2013 pairs
transcriptome and proteome on the same 22 sequenced isolates. McManus 2014 reads ribosome
occupancy beside mRNA abundance in \org{S. cerevisiae}, \org{S. paradoxus} and their F1
hybrid, so translation is measured allele by allele across two genomes; the
\org{S. paradoxus} arm is out of species and is stored apart or dropped. Foss 2007 is the
first segregant proteome on the BY x RM cross, fifteen years and two platforms before
Grossbach 2022, and its figures are the least confirmed of the eight: neither a segregant
count nor a proteome depth is stated in the abstract.

\subsection{What is still missing after the eight\secstatus{todo}}\label{sec:state-gaps}

Three things no row in the pool supplies. There is no proteome on a condition axis across
isolates; Teyssonniere 2024 and Muenzner 2024 are one medium each, and the only proteome
with a condition axis anywhere in the list, Lahtvee 2017, is wild type. There is no
translation layer on natural isolates beyond McManus 2014's three genomes. And there is no
single-cell proteome on any genome axis at all; the per-cell rows in the perturb-seq band
are transcriptomes. These are the shape of the 46-row shortfall, and discovery for the next
pass should start there.
